\documentclass[11pt,a4paper]{article}

\usepackage[margin=1in]{geometry}
\usepackage[T1]{fontenc}
\usepackage{lmodern}
\usepackage{booktabs}
\usepackage{tabularx}
\usepackage{array}
\usepackage{parskip}
\usepackage[hidelinks]{hyperref}

\newcolumntype{L}{>{\raggedright\arraybackslash}X}
\renewcommand{\arraystretch}{1.3}

\title{NewsRadar: Intermediate Project Update\\[4pt]
  \large A Geographical Visualization of News Around the Globe}
\author{Debasish Bordoloi\\CS 699 (Software Lab)}
\date{October 17, 2026}

\begin{document}

\maketitle

\section{Project Description}

NewsRadar is a CS 699 (Software Lab) team project: a web application that
scrapes news articles, extracts the places each article mentions, and plots
them on an interactive map. It is built as a three-stage Python pipeline
(scraping, location and topic mining, then a FastAPI backend with a
React + Leaflet map) around a shared PostgreSQL + PostGIS database. The
current version targets Indian news sources and geocodes places within India.

The main features of our application will be:

\begin{itemize}
    \item News articles collected from multiple sources will be mapped to the locations they discuss, allowing users to explore what is happening across different regions of the world.
    \item News events from around the world will be represented as interactive points (blips) on a global map. Selecting a point will allow users to view the relevant news articles associated with that location.
    \item News websites will be continuously scraped to collect and update news data from multiple sources.
    \item Multiple articles reporting on the same real-world event will be identified and clustered together, preventing duplicate representation of the same event on the map.
    \item News will be categorized into different categories, including but not limited to Business, Sports, Military, Disasters, Crime, Environment, Technology, and Opportunities.
    \item Users will be able to filter news based on time and date ranges, allowing them to view recent or historical events.
    \item Users will be able to zoom into any geographical region and explore the news and events occurring in that area.
    \item The platform will provide links to the original news articles, allowing users to access the complete information from the respective news sources.
    \item Users will be able to search for news using text-based queries. The system will identify relevant news articles from across the world and display the corresponding events geographically on the map. For example, a search for "earthquake" would show relevant earthquake-related news from different locations around the world.
\end{itemize}

\section{My Contributions}

I worked on four areas, from the map frontend down to geocoding.

\begin{center}
\begin{tabularx}{\textwidth}{@{}p{0.2\textwidth}L@{}}
\toprule
\textbf{Area} & \textbf{What I did} \\
\midrule
Frontend &
  Created the \texttt{frontend/} app from the React + Vite template with
  ESLint, added \texttt{leaflet} and \texttt{react-leaflet}, and wrote
  \texttt{App.jsx}: a full-screen map with OpenStreetMap tiles, scroll-wheel
  zoom and a minimum zoom level. Article markers will be drawn on this base. \\
Infrastructure &
  Made \texttt{docker-compose.yml} work with Podman by using the fully
  qualified image name \texttt{docker.io/postgis/postgis:16-3.4}. Podman does
  not default to Docker Hub, so the short name failed; Docker users are
  unaffected. \\
Database &
  Designed the tables that store locations and link them to articles
  (Section~\ref{sec:schema}). \\
Geocoding &
  Wrote \texttt{geocoding/geocode.py}, which turns place names into
  coordinates and caches them in the database (Section~\ref{sec:geocoding}). \\
\bottomrule
\end{tabularx}
\end{center}

\subsection{Location Schema}
\label{sec:schema}

I added two tables and three indexes to \texttt{db/init/01\_schema.sql}, so
each place is stored once and linked to every article that mentions it.

\begin{center}
\begin{tabularx}{\textwidth}{@{}p{0.36\textwidth}L@{}}
\toprule
\textbf{Object} & \textbf{Purpose} \\
\midrule
\texttt{locations} table &
  One row per distinct place: \texttt{location\_name}, \texttt{latitude},
  \texttt{longitude}. CHECK constraints keep coordinates in valid ranges. \\
\texttt{article\_location\_map} table &
  Many-to-many link between articles and locations.
  \texttt{mention\_order} records which place an article is mainly about.
  Deleting an article deletes its links. \\
Index on \texttt{lower(location\_name)} &
  Case-insensitive name lookup, which lets the table act as the geocoding
  cache. \\
GiST index on \texttt{point(longitude, latitude)} &
  Bounding-box queries for the current map view, and nearest-first queries. \\
Index on \texttt{article\_location\_map}\allowbreak\texttt{(location\_id)} &
  Fetching all articles about one place. \\
\bottomrule
\end{tabularx}
\end{center}

\subsection{Geocoding Module}
\label{sec:geocoding}

\texttt{geocoding/geocode.py} turns place names into coordinates using
OpenStreetMap's Nominatim service through \texttt{geopy}, with the
\texttt{locations} table acting as a cache.

\texttt{fetch\_coordinates(name, hint)} searches within India only. The
optional hint, such as a state, picks between places with the same name: for
example, ``Aurangabad'' with the hint ``Bihar'' returns the town in Bihar, not
the one in Maharashtra.

\texttt{get\_or\_create\_location(conn, name, hint)} checks the
\texttt{locations} table first and calls Nominatim only for names it has not
seen. New places are inserted and their \texttt{id} is returned, ready to link
through \texttt{article\_location\_map}. Names Nominatim cannot find are
stored with NULL coordinates so they are not looked up again, while network
errors are raised instead of stored.

The module follows Nominatim's usage policy: at most 1 request per second, an
identifying User-Agent, and cached results. It also runs from the command
line:

\begin{center}
\texttt{python geocoding/geocode.py Guwahati Majuli -{}-hint Assam}
\end{center}

\textbf{Testing.} I checked the lookups against the live Nominatim service with
Guwahati, Aurangabad with the Bihar hint, and a made-up name. The full insert
into the database has not been run yet.

\end{document}